\section*{3 A combinatorial lemma (author source pp. 13:3--13:4)}
Both Theorems 1 and 2 rely on a simple combinatorial lemma.
\begin{originalresult}{Lemma 3}
Let $\mathcal A\subseteq 2^{[n]}$ be a family of subsets of $[n]$. Assume that $1\le k\le n$ is such that every $k$-element subset of $[n]$ is contained in some $A\in\mathcal A$. Then there exists a subfamily $\mathcal A'\subseteq\mathcal A$ of size $|\mathcal A'|\le O(|\mathcal A|^{1/k}\log(n/k))$ with $\bigcup\mathcal A'=[n]$. In particular, if $|\mathcal A|\le 2^k$ then $|\mathcal A'|\le O(\log n)$.
\end{originalresult}
\begin{proof}
Assume that $|\mathcal A|\le a^k$. Let $t$ be the size of a largest set in $\mathcal A$. Then we have
\[\binom nk\le a^k\binom tk.\]
Hence $t\ge\frac n{ea}$, using the estimates $(\frac nk)^k\le\binom nk$, $\binom tk\le(\frac{et}k)^k$. Take some $A_0\in\mathcal A$ of size $t$. Let
\[\mathcal A_1:=\{A\setminus A_0:A\in\mathcal A\}.\]
Then every subset of $U_1:=[n]\setminus A_0$ of size at most $k$ is contained in some member of $\mathcal A_1$. The size of $U_1$ is at most $n(1-\frac1{ea})$. Similarly, take a largest set $A_1$ from $\mathcal A_1$ and obtain a new family $\mathcal A_2\subseteq 2^{U_2}$ on $U_2:=U_1\setminus A_1$. After $s$ steps, the size of $U_s$ is at most $n(1-\frac1{ea})^s$ and after $s\le O(a\log(n/k))$ steps we have $|U_s|\le k$. This guarantees that the largest set in $\mathcal A_s$ is $U_s$ itself and $[n]=\bigcup_{i=0}^s A_s$. By construction, every $A_i$ is contained in some element of the original family $\mathcal A$.
\end{proof}
